\section{Preliminaries}\label{sec:prelim}

This section fixes the notation, the codes, and the two maps from which the test is built.

\subsection{Notation}\label{sec:notation}

Every symbol of the paper is defined in this section or at its first use; \cref{app:notation}
lists them all.

\begin{itemize}
\item \emph{Numbers.} $\N:=\{0,1,2,\dots\}$, $\Z$ is the ring of integers, $[n]:=\{1,\dots,n\}$
for $n\in\N$, and $\frac1a\Z:=\{z/a:\ z\in\Z\}$ for an integer $a\ge1$. We set
$\max\emptyset:=-\infty$.
\item \emph{Fields.} $\F$ is a finite field and $\F^*:=\F\setminus\{0\}$ its multiplicative
group; $\E$ denotes a finite extension field of $\F$ when one occurs.
\item \emph{Multi-indices.} For $0\le i\le m$, a multi-index at level $i$ is a tuple
$\alpha=(\alpha_{i+1},\dots,\alpha_m)\in\N^{m-i}$, and
$x^\alpha:=x_{i+1}^{\alpha_{i+1}}\cdots x_m^{\alpha_m}$; at level $0$,
$\alpha=(\alpha_1,\dots,\alpha_m)$. For a \emph{weight vector} $w=(w_1,\dots,w_m)\in\N^m$,
$\langle w,\alpha\rangle:=\sum_{l=i+1}^mw_l\alpha_l$. The vector $e_l\in\N^m$ has $1$ in
coordinate $l$ and $0$ elsewhere, and $\mathbf 1:=(1,\dots,1)\in\N^m$.
\item \emph{Coefficients.} Let $R$ be a commutative $\F$-algebra. Every
$q\in R[x_{i+1},\dots,x_m]$ has a unique expansion $q=\sum_\alpha q_\alpha x^\alpha$ with
$q_\alpha\in R$, finitely many of them nonzero; $q_\alpha$ is the \emph{coefficient} of
$x^\alpha$, and $\deg_{x_l}q:=\max\{\alpha_l:\ q_\alpha\ne0\}$.
\item \emph{Univariate formats.} For $e\in\N$ and $t\ge1$,
$\F[T]_{\le e}:=\{g\in\F[T]:\deg g\le e\}$ and $\F[T]_{<t}:=\{g\in\F[T]:\deg g<t\}$, with
$\deg0:=-\infty$.
\item \emph{Words.} A word on a finite set $L$ with values in a set $A$ is a map $u:L\to A$;
for a polynomial $P$, $P|_L$ is the word $x\mapsto P(x)$.
\item \emph{Weight functions.} A weight function on $L$ is a map $\mu:L\to[0,1]$; it is not
to be confused with a weight vector $w$. For $A\subseteq L$,
$\mu(A):=\frac1{|L|}\sum_{x\in A}\mu(x)$. The weight function $\mu$ has \emph{common
denominator} $M\in\N\setminus\{0\}$ if $\mu(x)\in\frac1M\Z$ for every $x\in L$.
\end{itemize}

\begin{definition}[Polynomials]\label{def:poly}
$R$ denotes a commutative $\F$-algebra (in practice $R=\F$ or $R=\F[T]$, the ring of
polynomials in the formal variable $T$ that the folds introduce). For $1\le i\le m+1$,
$R[x_i,\dots,x_m]_{<k}:=\{q\in R[x_i,\dots,x_m]:\ \deg_{x_l}q<k\text{ for }i\le l\le m\}$; for
$i=m+1$ this is $R$.
\end{definition}

\subsection{Smooth domains}\label{sec:domains}

The test runs on a chain of multiplicative cosets whose lengths are divided by the fold arity
$k$ at each step. Smoothness makes every member of the chain a coset of a cyclic group, so a
point of $L_i$ has exactly $k$ preimages in $L_{i-1}$.

\begin{definition}[Field and domains]\label{def:domains}
$\F$ is a finite field; $k\ge 2$, $m\ge 1$ and $N$ are integers with $k^m\mid N$ and
$N\mid |\F|-1$. $L_0\subset\F^*$ is a coset of the subgroup of order $N$, and
$L_i:=\{x^k: x\in L_{i-1}\}$ for $1\le i\le m$, so $|L_i|=N/k^i$. The \emph{fibre} of
$y\in L_i$ is $\fib(y):=\{x\in L_{i-1}: x^k=y\}$; it equals $x_0\Omega_k$ for any
$x_0\in\fib(y)$, where $\Omega_k:=\{\zeta\in\F:\ \zeta^k=1\}$ is the group of $k$-th roots of unity, and $|\fib(y)|=k$.
\end{definition}

Every fibre is a coset of $\Omega_k$, because $x^k=y$ and $x_0^k=y$ give $(x/x_0)^k=1$. The
condition $k\mid N\mid|\F|-1$ also makes $k$ invertible in $\F$. Both facts are used in
\cref{lem:fold-sub}(a), where the components of Definition~\ref{def:fold} recover the
coefficients of a polynomial.

\subsection{Codes and agreement}\label{sec:codes}

Both codes are evaluation codes on a set $L\subseteq\F$: the Reed--Solomon code of degree
below $d$, and its interleaving, whose symbols are $t$-tuples. Values of oracle words
are sometimes polynomials in $T$ of degree below $t$; we read such a value as a vector of
$t$ coefficients.

\begin{definition}[Codes]\label{def:codes}
$\RS[\F,L,d]:=\{P|_L:\deg P<d\}$. The $t$-interleaved code $\RS[\F,L,d]^t$ is the set
of words $L\to\F^t$ whose $t$ coordinates lie in $\RS[\F,L,d]$; we identify
$\F[T]_{<t}$ with $\F^t$ through coefficients.
\end{definition}

Distance is relative Hamming distance: for words $u,v$ on a finite set $L$,
$\Delta(u,v):=\frac1{|L|}|\{x\in L:u(x)\ne v(x)\}|$, and for a set $C$ of words on $L$,
$\Delta(u,C):=\min_{v\in C}\Delta(u,v)$. The analysis measures closeness with weights, and the
definition below covers the unweighted case too.

{\sloppy
\begin{definition}[Weighted agreement]\label{def:agree}
For words $u,v$ on $L$ and $\mu:L\to[0,1]$,
$\agree_\mu(u,v):=\frac1{|L|}\sum_{x:\,u(x)=v(x)}\mu(x)$ (as in BCIKS20, \S7). For a set
$C$ of words, $\agree_\mu(u,C):=\max_{v\in C}\agree_\mu(u,v)$. Without a subscript,
$\mu\equiv1$: $\agree(u,v)=1-\Delta(u,v)$.
\end{definition}\par}

\subsection{The Kronecker map and the fold}\label{sec:kron-prelim}

A polynomial with bounded individual degrees is committed through one univariate polynomial,
and a fold of that polynomial is the image of a polynomial with one fewer variable. The
first map is a bijection because the exponents of a monomial are the base-$k$ digits of its
image's exponent.

\begin{definition}[Kronecker map]\label{def:kron}
For $0\le i\le m$ and $q\in R[x_{i+1},\dots,x_m]_{<k}$,
\[
U^{(i)}_q(Y):=q\big(Y,Y^{k},\dots,Y^{k^{m-i-1}}\big)\in R[Y],\qquad U_f:=U^{(0)}_f .
\]
For $i=m$, $q$ is a constant and $U^{(m)}_q=q$.
\end{definition}

The fold acts on a word on $L_{i-1}$ by averaging its values on a fibre against the powers
of one challenge. The challenge is an element of $R$, so it may be a polynomial in $T$.

\begin{definition}[Fold]\label{def:fold}
For $u:L_{i-1}\to R$ and $y\in L_i$, the \emph{components} of $u$ are
\[
u_j(y):=\frac1k\sum_{x\in\fib(y)}x^{-j}\,u(x)\qquad(0\le j<k),
\]
and for $c\in R$ the \emph{fold} is $\Fold_c(u):L_i\to R$,
$\Fold_c(u)(y):=\sum_{j<k}c^j\,u_j(y)$.
\end{definition}

For $k=2$ the two roots in $\fib(y)$ are $x$ and $-x$, so Definition~\ref{def:fold} reads
\[
\Fold_c(u)(y)=\frac{u(x)+u(-x)}{2}+c\,\frac{u(x)-u(-x)}{2x},
\]
the value being independent of which of the two points is called $x$; this is the fold of
FRI \cite{BBHR18}.

\subsection{Weighted degree and the degree code}\label{sec:wdeg}

Weighting the exponents of a monomial by $w$ interpolates between the total degree and the
individual degrees, and the code of the test consists of the Kronecker images whose weighted
degree is at most $D$. Both degree notions are needed: the weighted one bounds the test, and
the graded one bounds what an intermediate round may send.

\begin{definition}[Weighted and graded degree]\label{def:wdeg}
Let $w\in\N^m$. The \emph{weighted degree} of $f=\sum_\alpha f_\alpha x^\alpha\in\F[x_1,\dots,x_m]$ is
\[
\deg_w f:=\max\{\langle w,\alpha\rangle:\ f_\alpha\ne0\}.
\]
Let $0\le i\le m$. Every $q\in\F[T][x_{i+1},\dots,x_m]$ is uniquely
$q=\sum_{a\in\N}\sum_\alpha q_{a,\alpha}T^ax^\alpha$ with $q_{a,\alpha}\in\F$, and its
\emph{graded degree} is
\[
\gdeg_w q:=\max\{a+\langle w,\alpha\rangle:\ q_{a,\alpha}\ne0\}.
\]
Both equal $-\infty$ on the zero polynomial.
\end{definition}

\begin{definition}[Graded codes]\label{def:graded}
The \emph{graded codes} are
\[
\Gc_i^{w,D}:=\big\{\,U^{(i)}_q\big|_{L_i}\ :\ q\in\F[T][x_{i+1},\dots,x_m]_{<k},\ \gdeg_w q\le D\,\big\},
\]
and the \emph{degree code} $C_{w,D}$ is the set of $T$-free elements of $\Gc_0^{w,D}$.
\end{definition}

With $w=\mathbf 1$ the weighted degree is the total degree, and with $w=e_l$ it is the
degree $\deg_{x_l}f$. Every word of $\Gc_0^{w,D}$ is the image of a polynomial with
$\deg_{x_l}<k$ for all $l$, so $C_{w,D}\subseteq\RS[\F,L_0,k^m]$ whatever $w$ and $D$ are.

\subsection{Interactive oracle proofs of proximity}\label{sec:iop}

An interactive oracle proof of proximity (IOPP) for a code $C$ is a public-coin protocol
between a prover and a verifier in which each message of the prover is an oracle: the
verifier may read single entries of it, at a cost charged per entry, and the prover may be
unbounded. Completeness asks that an honest prover make the verifier accept with probability
$1$ when the input word lies in $C$. Soundness asks for an error $\varepsilon(\delta)$: no
prover may make the verifier accept with probability above $\varepsilon(\delta)$ when the input
word is $\delta$-far from $C$. The number of entries read is the query complexity,
the number of exchanges is the round complexity, and the model is the one of
\cite{BCS16,RRR16}.

The test of \cref{sec:test} is an IOPP for $C_{w,D}$. Its prover sends the words
$u_1,\dots,u_m$ together with the terminal polynomial $g$, its verifier sends the challenges
$r_1,\dots,r_m$, and each repetition of its query phase reads $k$ entries of each of
$u_0,\dots,u_{m-1}$ (one fibre per level) and compares the last fold with $g$. FRI \cite{BBHR18} is an IOPP for
Reed--Solomon codes with logarithmic round and query complexity, and the fold it uses is the
case $k=2$ of Definition~\ref{def:fold}.
